\clearpage\thispagestyle{empty}
\pdfbookmark[0]{Quellen zu Komet (5. und 12. Juli 1813)}{comet1813jul-sources}
\begin{center}\textbf{Quellen zu Komet (5. und 12. Juli 1813)}\end{center}
\noindent Diese KI-Übertragung ins heutige Deutsch folgt dem Auszug mit der Überschrift Comet in den \href{https://gdz.sub.uni-goettingen.de/id/PPN236007467}{\textit{Werken}, Band~VI} von Gauß (Göttingen, 1874), S.~368. Die Briefe vom 5. und 12. Juli 1813 erschienen zusammen zuerst als Nr.~IX der Zeitschrift \textit{Monatliche Correspondenz zur Beförderung der Erd- und Himmelskunde}, Band~XXVIII (Juli 1813), S.~\href{https://api.digitale-sammlungen.de/iiif/image/v2/bsb10538617_00103/full/1600,/0/default.jpg}{97}, \href{https://api.digitale-sammlungen.de/iiif/image/v2/bsb10538617_00104/full/1600,/0/default.jpg}{98}, \href{https://api.digitale-sammlungen.de/iiif/image/v2/bsb10538617_00105/full/1600,/0/default.jpg}{99}. Die Links führen zum digitalisierten Exemplar der Bayerischen Staatsbibliothek.
\par\smallskip\noindent Die Werkausgabe kennzeichnet Auslassungen durch Striche. Der Erstdruck enthält auf S.~97--98 Ponds Sterntabelle und die zugehörige Erläuterung. Die Beobachtungen und Elemente des Kometen stehen auf S.~99. Die beiden Einträge zum 25. April sind in beiden Ausgaben verschieden. Das Kometenzeichen folgt der Werkausgabe. KI-Übertragung und KI-Satz.
\clearpage
